\subsection{Euler-Poincaré characteristics}

In this n°, we consider a set C of classes of A-modules that is additive and left exact, that is, satisfying the following two conditions:

\begin{enumerate}
\def\labelenumi{(\Alph{enumi})}
\item
  If M and N are two A-modules of type C, then M ⊕ N is of type C.
\item
  If \(0 \to \mathbf{M}' \to \mathbf{M} \to \mathbf{M}'' \to 0\) is an exact sequence of A-modules and if \(\mathbf{M}\) and \(\mathbf{M}''\) are of type \(\mathcal{C}\) , then \(\mathbf{M}'\) is of type \(\mathcal{C}\) .
\end{enumerate}

We say that \(\mathcal{C}\) is stable if it satisfies the following conditions, which imply (A) and (G):

\begin{enumerate}
\def\labelenumi{(\Alph{enumi})}
\setcounter{enumi}{4}
\item
  ("C is stable under extensions.") If 0 → M' → M → M'' → 0 is an exact sequence of A-modules and if M' and M'' are of type C, then M is of type C.
\item
  ("C is stable under kernels and cokernels.") For every homomorphism f of A-modules of type C, the A-modules Ker f and Coker f are of type C.
\end{enumerate}

We denote by K( \(\mathcal{C}\) ) the Grothendieck group of \(\mathcal{C}\) and {[}M{]} \(_{\mathcal{C}}\) or {[}M{]} the element of K( \(\mathcal{C}\) ) defined by the A-module M (VIII, § 6, no. 2). Let G be a commutative group and \(\varphi\) a homomorphism from K( \(\mathcal{C}\) ) into G.

Examples. --- 1) If A is a field, one may take for C the set of classes of finite-dimensional vector spaces, and for φ the isomorphism of K(C) onto Z defined by φ({[}M{]}) = dim (M).

\begin{enumerate}
\def\labelenumi{\arabic{enumi})}
\setcounter{enumi}{1}
\tightlist
\item
  One may take for \(\mathcal{C}\) the set of classes of modules of finite length, and for \(\varphi: \mathbf{K}(\mathcal{C}) \to \mathbf{Z}\) the homomorphism defined by \(\varphi([\mathbf{M}]) = \text{long}_{\mathbf{A}}(\mathbf{M})\) .
\end{enumerate}

One says that a graded A-module M is of type C if \(M_{n}\) is of type C for every n (for this, it is necessary, if M is bounded, and sufficient, if C is stable, that the module M be of type C).

DEFINITION 8. --- Let M be a bounded graded A-module of type C and let (M \(_{n}\) ) be its grading. One calls the φ-characteristic of M, and denotes by χφ(M) or simply χ(M), the element ∑ (-1) \(^{n}\) φ({[}M \(_{n}\) {]}) of G.

This definition applies in particular when M is the graded module underlying a complex of A-modules.

Examples. --- 3) If M is bounded of type C, the same is true of M(p) for every \(p \in Z\) , and we have \(\chi(\mathbf{M}(p)) = (-1)^{p} \chi(\mathbf{M})\) .

\begin{enumerate}
\def\labelenumi{\arabic{enumi})}
\setcounter{enumi}{3}
\tightlist
\item
  Let \(0 \to \mathbf{M}' \to \mathbf{M} \to \mathbf{M}'' \to 0\) be an exact sequence of graded A-modules and graded homomorphisms of degree 0. If \(\mathbf{M}, \mathbf{M}'\) and \(\mathbf{M}''\) are bounded of type \(\mathcal{C}\) , we have
\end{enumerate}

\[
\chi (\mathbf {M}) = \chi (\mathbf {M} ^ {\prime}) + \chi (\mathbf {M} ^ {\prime \prime}).
\]

If M and M'' are bounded of type C, the same is true of M'; if C is stable and two of the three modules are bounded of type C, the same is true of the third.

\begin{enumerate}
\def\labelenumi{\arabic{enumi})}
\setcounter{enumi}{4}
\tightlist
\item
  Let \(u: \mathbf{C}' \to \mathbf{C}\) be a morphism of bounded complexes of type \(\mathcal{C}\) . Then Con \((u)\) is bounded of type \(\mathcal{C}\) , and we have:
\end{enumerate}

\[
\chi (\operatorname{Con} (u)) = \chi (\mathbf {C}) - \chi (\mathbf {C} ^ {\prime}).
\]

\begin{enumerate}
\def\labelenumi{\arabic{enumi})}
\setcounter{enumi}{5}
\tightlist
\item
  One can take for G the group K(C) itself, and for φ the identity; we then denote by \(\chi_{\mathcal{C}}(\mathrm{M})\) the element \(\chi_{\varphi}(\mathrm{M}) = \sum (-1)^{n}[\mathrm{M}_{n}]\) of K(C).
\end{enumerate}

Remark. --- One calls the Poincaré polynomial of M relative to \(\varphi\) the element \(\mathrm{P}_{\mathbf{M}}(t) = \sum \varphi([\mathrm{M}_{n}]) t^{n} \in \mathrm{G} \otimes \mathbf{Z}[t, t^{-1}]\) . We have \(\mathrm{P}_{\mathbf{M}}(1) = \varphi([\mathrm{M}])\) and \(\mathrm{P}_{\mathbf{M}}(-1) = \chi(\mathrm{M})\) .

Lemma 4. --- Let C be a bounded complex of type \(\mathcal{C}\) . If H(C) = 0, we have \(\chi(\mathbf{C}) = 0\) . This follows from VIII, § 6, n° 1, cor. of prop. 1.

PROPOSITION 10. --- Let C and C' be two bounded complexes of type C. If there exists a homologism \(u: \mathbf{C}' \to \mathbf{C}\) , we have \(\chi(\mathbf{C}) = \chi(\mathbf{C}')\) .

Indeed, Con \((u)\) is bounded of type \(\mathcal{C}\) and we have \(\chi(\text{Con } (u)) = \chi(\text{C}) - \chi(\text{C}')\) ; on the other hand, H(Con \((u)) = 0\) by X, p.~38, cor., hence \(\chi(\text{Con } (u)) = 0\) (lemma 4).

PROPOSITION 11. --- Let C be a bounded complex of type C.

\begin{enumerate}
\def\labelenumi{\alph{enumi})}
\item
  If \(\mathcal{C}\) is stable, H(C) is of type \(\mathcal{C}\) .
\item
  If H(C) is of type C, then so are B(C) and Z(C), and we have \(\chi(\mathrm{H}(\mathrm{C})) = \chi(\mathrm{C})\) .
\item
  If \(\mathcal{C}\) is stable, for every \(n\) the module \(\mathbf{Z}_n(\mathbf{C})\) is of type \(\mathcal{C}\) as kernel of \(d_n: \mathbf{C}_n \to \mathbf{C}_{n-1}\) , and \(\mathrm{H}_n(\mathbf{C})\) is of type \(\mathcal{C}\) as cokernel of \(\mathbf{C}_{n+1} \to \mathbf{Z}_n\) . On the other hand, \(\mathrm{H}_n(\mathbf{C}) = 0\) as soon as \(\mathbf{C}_n = 0\) .
\item
  Suppose H(C) is of type C. The canonical exact sequences:
\end{enumerate}

\[
0 \to Z _ {n} (\mathbf {C}) \to C _ {n} \to B _ {n - 1} (\mathbf {C}) \to 0
\]

\[
0 \to \mathrm{B} _ {n} (\mathrm{C}) \to \mathrm{Z} _ {n} (\mathrm{C}) \to \mathrm{H} _ {n} (\mathrm{C}) \to 0
\]

show by induction on n starting from the right bound of C that \(Z_{n}(C)\) and \(B_{n}(C)\) are of type C for every n.~We then have

\[
\chi (\mathrm{C}) = \chi (\mathrm{Z} (\mathrm{C})) + \chi (\mathrm{B} (\mathrm{C}) (- 1)) = \chi (\mathrm{Z} (\mathrm{C})) - \chi (\mathrm{B} (\mathrm{C})) = \chi (\mathrm{H} (\mathrm{C})).
\]

COROLLARY. --- If C is stable and C is bounded of type C, the graded module H(C) is bounded of type C and we have \(\chi(\mathrm{H}(\mathrm{C})) = \chi(\mathrm{C})\) .

PROPOSITION 12. --- Let \(0 \to C' \to C \to C'' \to 0\) be an exact sequence of complexes. a) If H(C), H(C') and H(C'') are bounded of type C, we have

\[
\chi (\mathrm{H} (\mathbf {C})) = \chi (\mathrm{H} (\mathbf {C} ^ {\prime})) + \chi (\mathrm{H} (\mathbf {C} ^ {\prime \prime})) .
\]

\begin{enumerate}
\def\labelenumi{\alph{enumi})}
\setcounter{enumi}{1}
\tightlist
\item
  If C is stable, and if two of the graded modules H(C), H(C') and H(C'') are bounded of type C, then so is the third.
\end{enumerate}

Part a) follows from lemma 4 applied to the zero-homology complex defined by the homology exact sequence associated with the given exact sequence. Part b) follows, by considering this homology exact sequence, from the following lemma:

Lemma 5. --- Let M → N → P → Q → R be an exact sequence of A-modules. If C is stable, and if M, N, Q and R are of type C, then the module P is of type C.

Set \(\mathbf{N}' = \operatorname{Coker} (\mathbf{M} \to \mathbf{N})\) and \(\mathbf{Q}' = \operatorname{Ker} (\mathbf{Q} \to \mathbf{R})\) . The modules \(\mathbf{N}'\) and \(\mathbf{Q}'\) are of type \(\mathcal{C}\) , and we have an exact sequence \(0 \to \mathbf{N}' \to \mathbf{P} \to \mathbf{Q}' \to 0\) .

COROLLARY. --- Suppose C is stable, and let \(u : C' \to C\) be a morphism of complexes such that \(\mathrm{H}(\mathbf{C})\) and \(\mathrm{H}(\mathbf{C}')\) are bounded of type C. Then \(\mathrm{H}(\mathrm{Con}(u))\) is bounded of type C, and we have

\[
\chi (\mathrm{H} (\text { Con } (u))) = \chi (\mathrm{H} (\mathrm{C})) - \chi (\mathrm{H} (\mathrm{C} ^ {\prime})) .
\]

This follows from prop. 12 applied to the exact sequence of complexes (X, p.~37, prop. 7)

\[
0 \rightarrow \mathrm{C} \rightarrow \operatorname{Con} (u) \rightarrow \mathrm{C} ^ {\prime} (- 1) \rightarrow 0.
\]

Remark. --- Let E be a complex, \(h: \mathrm{E} \to \mathrm{C}\) and \(h': \mathrm{E} \to \mathrm{C}'\) homotopisms with C and C' bounded of type \(\mathcal{C}\) . We then have \(\chi(\mathrm{C}) = \chi(\mathrm{C}')\) . Indeed, if \(h_1\) is an inverse of \(h\) up to homotopy, \(h' \circ h_1\) is a homotopism, hence a homologism from C to

C' and one can apply prop. 10. Consequently, one can extend definition 8 by setting \(\chi(E) = \chi(C)\) as soon as there exists a homotopism from E to a bounded complex C of type C. Propositions 10, 11, 12 and their corollaries generalize to this setting.

Application:

* Let X be a topological space admitting a finite cell decomposition (cf.~n° 3). a) Let K and K' be two fields, set \(b_i = \dim_K(\mathrm{H}_i(\mathrm{X}, \mathrm{K}))\) and \(b_i' = \dim_K(\mathrm{H}_i(\mathrm{X}, \mathrm{K}'))\) . One does not necessarily have \(b_i = b_i'\) , but one has \(\Sigma(-1)^i b_i = \Sigma(-1)^i b_i'\) .

\begin{enumerate}
\def\labelenumi{\alph{enumi})}
\setcounter{enumi}{1}
\tightlist
\item
  Let \((\mathbf{X}_{n})\) and \((\mathbf{X}_{n}^{\prime})\) be two finite cell decompositions of X, and let \(c_{n}\) and \(c_{n}^{\prime}\) be the number of cells of dimension n in these two decompositions. We have
\end{enumerate}

\[
\Sigma (- 1) ^ {i} c _ {i} = \Sigma (- 1) ^ {i} c _ {i} ^ {\prime}.
\]

\begin{enumerate}
\def\labelenumi{\alph{enumi})}
\setcounter{enumi}{2}
\tightlist
\item
  With the notations of \(a)\) and \(b)\) , we have \(\Sigma(-1)^{i} c_{i} = \Sigma(-1)^{i} b_{i}\) .
\end{enumerate}

The properties \(a\) ) and \(b\) ) follow from \(c\) ), and \(c\) ) follows from prop. 11 applied to the complex \(\Gamma\) described in no. 3, taking for \(\mathscr{C}\) the class of finite-dimensional K-vector spaces and for \(\varphi\) the function defined by \(\varphi([\mathrm{M}]) = \dim_{\mathrm{K}}(\mathrm{M})\) (X, p.~40, example 1). *

